\documentclass[english,twoside]{article}
\usepackage[utf8]{inputenc}
\usepackage[a4paper, left=3cm, right=3cm, top=4cm]{geometry}

% Language & links
\usepackage[english]{babel}
\usepackage[hidelinks]{hyperref}
\usepackage{etoolbox}

% Math
\usepackage{amsmath, amsfonts, amssymb, mathtools}
\usepackage{mathrsfs}      % script letters
\usepackage{physics}       % physics notation
\usepackage{tikz-cd}       % commutative diagrams
\usepackage{amsthm}        % proof

% Figures
\usepackage{graphicx} 
\usepackage{caption}
\usepackage{subcaption}
\usepackage[table,dvipsnames]{xcolor}
\usepackage{float}

% Tables
\usepackage{array}
\usepackage{multirow}
\usepackage{colortbl}
\usepackage{hhline}
\usepackage{booktabs}
\usepackage{booktabs}
\usepackage[table]{xcolor}

% Layout
\usepackage{tcolorbox}     % colored boxes
\usepackage{enumitem}      % custom lists
\usepackage{epigraph}      % quotes at section start
\usepackage{fancyhdr}      % 
\usepackage{setspace}      % interlinea 

% Code
\usepackage{minted}

% Appendix
\usepackage[toc]{appendix}

% Bibliography
\usepackage{csquotes}
\usepackage[
    backend=biber,
    style=alphabetic
]{biblatex}
\addbibresource{references.bib}

\theoremstyle{plain}
\newtheorem{theorem}{Theorem}[section]
\newtheorem{lemma}[theorem]{Lemma}
\newtheorem{proposition}[theorem]{Proposition}
\newtheorem{corollary}[theorem]{Corollary}

\theoremstyle{definition}
\newtheorem{definition}[theorem]{Definition}

\theoremstyle{remark}
\newtheorem{remark}[theorem]{Remark}


\numberwithin{equation}{section} % Numeration eq$

\onehalfspacing % Spaziatura

\setlength{\headheight}{13.6pt}

% Titolo, autore e data PRIMA di definire \thetitle
\title{Classical Integrable Systems}
\author{Tommaso Cleva\footnote{tommasocleva@studenti.units.it}}
\date{\today}

% Copia il contenuto di \title dentro \thetitle, per usarlo in fancyhdr
\makeatletter
\let\thetitle\@title
\makeatother

\pagestyle{fancy}
\fancyhead[LE]{\nouppercase{\leftmark}}   % sezione su pagine pari
\fancyhead[RO]{\nouppercase{\leftmark}}   % sezione su pagine dispari
\fancyhead[LO]{\textit{\thetitle}}        % titolo su pagine dispari
\fancyhead[RE]{\textit{\thetitle}}        % titolo su pagine pari
\fancyfoot[C]{\thepage}


\begin{document}

\begin{titlepage}
    \centering

    \vspace*{2cm}

    {\Huge \bfseries 
    Test
    }\\[1.5cm]

    {\large
    Tommaso Cleva\\[0.3cm]
    \textit{University of Trieste and SISSA}\\[0.3cm]
    \texttt{tommaso.cleva@studenti.units.it}
    }

    \vspace{2cm}

    {\large
    Supervisor: Mikhail Bershtein
    }

    
    \vfill 
    
    {\large \today}


\end{titlepage}

\newpage
\begin{abstract}
    In this thesis, we give an introduction to classical integrable systems. The topics include Liouville integrability, Lax representations, the $r$-matrix formalism, the classical and modified Yang--Baxter equations, and the open Toda chain.
\end{abstract}
\newpage

\tableofcontents
\newpage

\section{Introduction}
Integrable systems occupy a central role in classical and mathematical physics, providing rare examples of nonlinear dynamical systems whose evolution can be described exactly \cite{babelon2003introduction,perelomov1990integrable}. Their origin lies in classical and celestial mechanics, where the search for explicit solutions to the equations of motion led to the identification of systems admitting sufficiently many conserved quantities. This program culminated in the Liouville–Arnold theorem, which gives a geometric characterization of integrability in terms of invariant tori and action–angle variables \cite{babelon2003introduction}.

A decisive conceptual shift occurred in the second half of the twentieth century, when integrability was recognized not merely as solvability by quadratures, but as the manifestation of an underlying algebraic structure. A key development in this direction is the introduction of Lax pairs, which encode the dynamics in a matrix equation of the form $\dot{L} = [M,L]$ \cite{https://doi.org/10.1002/cpa.3160210503}. This formulation, originally introduced in the study of nonlinear evolution equations such as the Korteweg–de Vries equation \cite{gardner1967method,Zakharov1971}, allows one to interpret the time evolution as an isospectral deformation and to construct conserved quantities from the spectral invariants of the Lax matrix.

This algebraic viewpoint is further refined by the $r$-matrix formalism, which organizes the Poisson brackets of the entries of the Lax matrix into a tensorial structure \cite{SemenovTyanShanskii1983,Sklyanin1982}. In this framework, the involutivity of the spectral invariants follows from a precise algebraic relation involving the $r$-matrix. These structures reveal a deep connection between integrable systems and Lie algebras, with the tensor Casimir element and related invariant tensors playing a central role \cite{perelomov1990integrable}.

The aim of this work is to present these ideas in a concrete and largely self-contained manner, with particular emphasis on the interplay between Hamiltonian dynamics, Lax representations, and algebraic structures. After reviewing Liouville integrability and the basic properties of Lax pairs, we develop the $r$-matrix formalism, discuss the role of the Casimir element, and derive the classical and modified Yang–Baxter equations. These concepts are then illustrated on the example of the open Toda chain \cite{toda1967vibration,Olshanetsky1979,KOSTANT1979195}.

While Sections~\ref{subsec:Introduction}--\ref{subsec:Action--angle variables} and Section~\ref{sec:Commuting flows} follow standard presentations (mainly based on \cite{babelon2003introduction}), several parts of this work are developed independently. In particular, Section~\ref{subsec:Example: Lax pairs for oscillators} provides a self-contained derivation of a Lax representation for one-dimensional systems, while Sections~\ref{sec:rmatrix_formalism}, \ref{sec:casimir}, and \ref{sec:YB} present a direct and explicit treatment of the $r$-matrix formalism, the tensor Casimir element, and the classical and modified Yang–Baxter equations.

The proof of the $r$-matrix structure for the open Toda chain in Section~\ref{sec:toda_rmatrix} is given in detail, and Section~\ref{sec:Independence of Hk} contains a direct argument for the independence of the Hamiltonians. The goal is to provide a coherent and partially self-contained presentation, with emphasis on explicit constructions and derivations.

\paragraph{Acknowledgements}
I would like to express my deepest gratitude to my supervisor, Mikhail Bershtein, for his invaluable guidance, insightful discussions, and continuous support throughout this work. I am especially grateful for his exceptional kindness and generosity with his time, as well as his constant availability. His encouragement and openness have made this experience not only intellectually enriching, but also deeply rewarding on a personal level.

\section{Integrable dynamical systems}
\subsection{Introduction}\label{subsec:Introduction}
In Classical Mechanics, the state of a system is represented by a point in phase space, whose coordinates are the positions $q_i$ and the momenta $p_i$ \cite{babelon2003introduction,perelomov1990integrable}.  
The Hamiltonian $H(p_i, q_i)$ is a function on this space, and the equations of motion take the form:

\begin{equation} \label{eq:Hamilton}
    \dot{q_i} = \pdv{H}{p_i}, \qquad \dot{p_i} = -\pdv{H}{q_i}.
\end{equation}
A dot denotes a time derivative. For any function $F(p, q)$ depending on the phase-space variables, its time evolution is given by:
\begin{equation}
    \dot{F} = \dv{F}{t} = \{F,H\}.
\end{equation}
\begin{definition}
    The Poisson bracket between two functions $F$ and $G$ is defined as:
    \begin{equation}
        \{F, G\} = \sum_i \left(
        \pdv{F}{q_i} \pdv{G}{p_i} -
        \pdv{F}{p_i} \pdv{G}{q_i}
        \right).
    \end{equation}
\end{definition}
For the canonical coordinates we have:
\begin{equation} \label{eq:Foundamental_Poisson}
    \{q_i, q_j\} = 0, \qquad
    \{p_i, p_j\} = 0, \qquad
    \{q_i, p_j\} = \delta_{ij}.
\end{equation}
The Hamiltonian $H(p, q)$ is conserved during time evolution, since
\begin{equation}
    \dv{H}{t} = \{H,H\} = 0,
\end{equation}
so the motion remains confined to the subvariety $H = E$ in phase space.


\subsection{The Liouville Arnold theorem}

We consider a Hamiltonian system with phase space $M$ of dimension $2n$ \cite{babelon2003introduction,perelomov1990integrable}.  
Using canonical coordinates $p_i, q_i$, the non-degenerate Poisson bracket is defined as in equation (\ref{eq:Foundamental_Poisson}).  
A non-degenerate Poisson bracket on $M$ corresponds to a closed, non-degenerate 2-form $\omega$ (the symplectic form) satisfying $\dd \omega = 0$.  
In canonical coordinates it is given by
\begin{equation}
    \omega = \sum_j \dd p_j \wedge \dd q_j.
\end{equation}
\begin{definition}[Liouville integrability]
    Let $H$ be the Hamiltonian of the system. A system is \emph{Liouville integrable} if it admits $n$ independent conserved quantities $F_i$ $(i = 1, \dots, n)$ in involution
    \begin{equation}
        \{F_i, H\} = 0, \qquad \{F_i, F_j\} = 0.
    \end{equation}
\end{definition}

\begin{theorem}[The Liouville Arnold theorem]
     The solution of the equations of motion of a Liouville integrable system is obtained by "quadrature".
\end{theorem}

\begin{proof}
    Let 
    \begin{equation} \label{eq:def_alpha}
        \alpha = \sum_i p_i \dd{q_i}
    \end{equation}
    be the canonical 1-form and $\omega = \dd{\alpha} =\sum_i \dd{p_i} \wedge \dd{q_i}$ be the symplectic 2-form on the phase space $M$. Now we will construct a canonical transformation
    \begin{equation}
        (p_i,q_i) \mapsto (F_i, \psi_i),
    \end{equation}
    such that the conserved quantities $F_i$ are among the new coordinates:
    \begin{equation}
        \omega = \sum_i \dd{p_i} \wedge \dd{q_i} = \sum_i \dd{F_i} \wedge \dd{\psi_i}.
    \end{equation}
    If we succeed in doing that, the equations of motion become trivial:
    \begin{equation}
        \begin{cases}
            \dot{F}_j = \{F_j, H\} = 0 \\
            \dot{\psi}_j = \{\psi_j, H \} = \pdv{H}{F_j} = \Omega_j
        \end{cases}.
    \end{equation}
    The $\Omega_j$ depend only on $F$ and so are constant in time. In these coordinates, the solutions of the equation of motion are:
    \begin{equation}
        \begin{cases}
            F_j(t) = F_j(0) \\
            \psi_j(t) = \psi_j(0) + t \, \Omega_j
        \end{cases}.
    \end{equation}
    To construct this canonical transformation we exhibit the \textit{generating function} $S$. Let $M_f$ be the level manifold $F_i(p,q) = f_i$. Suppose that on $M_f$ we can solve for $p_i$, $p_i = p_i(f,q)$ and consider the function
    \begin{equation}
        S(F,q) \equiv \int_{m_0}^{m} \alpha = \int_{q_0}^{q} \sum_i p_i(f,q) \dd{q_i},
    \end{equation}
    where the integration path is drawn on $M_f$ and goes from the point of coordinate $(p(f,q_0),q_0)$ to the point $(p(f,q),q)$ where $q_0$ is some reference value.
    Assume that this function is well-defined, i.e. it is independent of the integration path from $m_0$ to $m$, then $p_j = \pdv{S}{q_j}$.
    Defining $\psi_i$ by
    \begin{equation}
        \psi_i = \pdv{S}{F_i},
    \end{equation}
    we have
    \begin{equation}
        \dd{S} = \sum_j \psi_j \dd{F_j} + p_j \dd{q_j}.
    \end{equation}
    Since $\dd^2 S = 0$ we deduce that $\omega = \sum_j \dd{p_j} \wedge \dd{q_j} = \sum_j \dd{F_j} \wedge \dd{\psi_j}$.
    This shows that if $S$ is a well-defined function, then the transformation is canonical.
    To show that $S$ exists, we must prove that it is independent of the integration path. By Stokes theorem, we have to prove that:
    \begin{equation}
        \dd{\alpha}|_{M_f} = \omega|_{M_f} = 0.
    \end{equation}
    Let $X_i$ be the Hamiltonian vector field associated with $F_i$, defined by $\omega(X_i, \cdot) = - \dd F_i$
    \begin{equation}
        X_i = \sum_k  \pdv{F_i}{p_k} \pdv{}{q_k} - \pdv{F_i}{q_k} \pdv{}{p_k}.
    \end{equation}
    These vector fields are tangent to the manifold $M_f$ because the $F_j$ are in involution,
    \begin{equation}
        X_i(F_j) = \{F_j, F_i\} = 0.
    \end{equation}
    Since the functions $F_1, \dots, F_n$ are independent, the differentials $\dd F_i$ are linearly independent on $M_f$. 
    By non-degeneracy of the symplectic form, this implies that the corresponding Hamiltonian vector fields $X_i$ are linearly independent and span the tangent space to the submanifold $M_f$ at each regular point $m \in M_f$. 
    But then $\omega(X_i, X_j) = \{F_i, F_j\} = 0$, so $\omega|_{M_f} = 0$, and therefore $S$ is locally well-defined.
\end{proof}
\begin{remark}
    From the closedness of $\alpha$ on $M_f$, the function $S$ is unchanged by \emph{continuous} deformations of the path $(m_0, m)$. However, if $M_f$ has non-trivial cycles, which is generically the case, $S$ is a multivalued function defined in a neighbourhood of $M_f$.
    The variation over a cycle
    \begin{equation}
        \Delta_{\text{cycle}} S = \int_{\text{cycle}} \alpha
    \end{equation}
    is a function of $F$ only. This induces a multivaluedness of the variables 
    $\psi_j$:
    \begin{equation}
        \Delta_{\text{cycle}} \psi_j = \frac{\partial}{\partial F_j} 
        \Delta_{\text{cycle}} S;
    \end{equation}
\end{remark}
\begin{remark}
    The definition we have given of a Liouville integrable system requires some care.
    Given any Hamiltonian $H$, the Darboux theorem implies that we
    can always find locally a system of canonical coordinates on phase space $(P_1, \dots, P_n; Q_1, \dots, Q_n)$, with $H = P_1$, hence fulfilling the hypothesis of the Liouville theorem. For integrable systems we require that the conserved quantities are globally defined on a sufficiently large open set, and that the surfaces $F_i = f_i$ are well-behaved and foliate the phase space.
    This is not generally the case for the $P_i$ constructed by the Darboux theorem. Moreover, in many standard examples, the conserved quantities are even algebraic functions of canonical coordinates on some open domain and the solutions of the equations of motion are analytic.
\end{remark}
\begin{remark}
    Using the Poisson commuting functions $ F_i $, one can solve \emph{simultaneously} the $ n $ time evolution equations $ \dv{F}{t_i} = \{F, F_i\} $, since:
    \begin{equation}
        \dv{}{t_i} \dv{}{t_j} F - \dv{}{t_j} \dv{}{t_i} F
        = \{\{F, F_j\}, F_i\} - \{\{F, F_i\}, F_j\}
        = \{F, \{F_j, F_i\}\} = 0.
    \end{equation}
    
    Since the Hamiltonian vector fields are well-defined and linearly independent on $M_f$, the flows define a locally free (no fixed points) and transitive (goes everywhere) action of a small open set in $ \mathbb{R}^n $ on the surface $ M_f $. Assuming that $ M_f $ is connected and compact, the flows extend to all values of the times $ t_i $ and fill the whole surface $ M_f $, hence we have a surjective action of $ \mathbb{R}^n $ on $ M_f $. The stabilizer of a point is an Abelian discrete subgroup of $ \mathbb{R}^n $ since the action is locally free, so it is of the form $ \mathbb{Z}^n $. Thus $ M_f $ appears as the quotient of $ \mathbb{R}^n $ by $ \mathbb{Z}^n $, i.e.\ a torus.
    
    This refinement, due to Arnold, of the Liouville theorem shows that, under suitable global hypothesis, the phase space is indeed foliated by $ n $-dimensional tori, called the Liouville tori. It is remarkable that for small perturbations of integrable systems, there still exist Liouville tori “almost everywhere”. This is the content of the famous Kolmogorov–Arnold–Moser (KAM) theorem.
\end{remark}
\subsection{Action--angle variables} \label{subsec:Action--angle variables}
As already noticed in the proof of the Liouville theorem, the level manifold $M_f$ has non-trivial cycles. Under suitable compactness and connectivity conditions, the $M_f$ are $n$-dimensional tori $T_n$. This suggests the introduction of angle variables to describe the motion along the cycles. The torus $T_n$ is isomorphic to a product of $n$ circles $C_i$. We may therefore choose special angular coordinates on $M_f$ that are dual to the $n$ fundamental cycles $C_i$ (see eq. (\ref{eq:angular_variables})).\\
The action variables $I_j$ are defined as the integrals of the canonical one-form over the cycles $C_j$:
\begin{equation}
    I_j = \frac{1}{2\pi} \int_{C_j} \alpha.
\end{equation}
The $I_j$ are functions of the constants of motion $F_j$, and we suppose they are independent, so that if the values of $I_j$ ($j = 1, \ldots, n$) are known, then the manifold $M_f$ is determined.\\
Let us consider the canonical transformation generated by the same function as above:
\begin{equation}
    S(I, q) = \int_{m_0}^{m} \alpha,
\end{equation}
but now expressed in terms of the variables $I_i$ instead of $F_i$.  
Denoting by $\theta_j$ the variable conjugate to $I_j$, the canonical transformation generated by $S$ is defined by
\begin{equation}
    p_k = \frac{\partial S}{\partial q_k}, 
    \qquad 
    \theta_k = \frac{\partial S}{\partial I_k}.
\end{equation}
The variables $\theta_k$ are canonically conjugate to the action variables $I_j$.  
We show that they can be regarded as normalized angular variables on the cycles $C_j$.  
That is,
\begin{equation}\label{eq:angular_variables}
    \frac{1}{2\pi} \int_{C_j} \dd{\theta_k} = \delta_{jk}.
\end{equation}
By definition of $\theta_k$, and assuming that the cycles $C_j$ can be chosen locally so as to depend smoothly on the action variables $I_k$, we can differentiate under the integral sign:
\begin{equation}
    \int_{C_j} \dd{\theta_k} 
    = \pdv{}{I_k} \int_{C_j} \dd{S}, 
    \qquad 
    \dd{S} = \sum_i \pdv{S}{q_i} \dd{q_i} + \pdv{S}{I_i} \dd{I_i}.
\end{equation}
Since on the manifold $M_f$ we have $dI_i = 0$, it follows that
\begin{equation}
    \int_{C_j} \dd{\theta_k}
    = \frac{\partial}{\partial I_k} 
    \int_{C_j} 
    \frac{\partial S}{\partial q_i} \dd{q_i} 
    = \frac{\partial}{\partial I_k} 
    \int_{C_j} \alpha 
    = 2\pi \delta_{jk}.
\end{equation}
This proves that the variables $\theta_k$ are indeed angle variables.
\subsubsection{Example: harmonic oscillator}
For the one-dimensional harmonic oscillator
\begin{equation}
    H = \frac{p^2}{2m} + \frac{m}{2}\omega^2 q^2 = E,
\end{equation}
the level sets of the Hamiltonian in the phase space $(q,p)$ are ellipses of the form
\begin{equation}
    \frac{q^2}{A^2} + \frac{p^2}{(m\omega A)^2} = 1,
\end{equation}
where $A$ is the amplitude, so that
\begin{equation}
    E = \frac{1}{2} m \omega^2 A^2.
\end{equation}
Since the motion is periodic, the level manifold is a one-dimensional torus (a circle), and there is a single action variable:
\begin{equation}
    I = \frac{1}{2\pi} \oint_C \alpha = \frac{1}{2\pi} \oint_C p \, \dd q.
\end{equation}
In one degree of freedom, the integral $\oint_C p \, \dd q$ represents the oriented area enclosed by the orbit in the $(q,p)$ plane. Therefore
\begin{equation}
    \oint_C p \, \dd q = \pi A \cdot m\omega A = \frac{2\pi E}{\omega},
\end{equation}
and hence
\begin{equation}
    I = \frac{E}{\omega}.
\end{equation}
The corresponding angle variable evolves linearly in time:
\begin{equation}
    \theta = \omega t + \theta_0.
\end{equation}

\subsection{Lax Pairs}

In many integrable systems, the time evolution can be described using a pair of matrices or operators $(L, M)$, known as a \emph{Lax pair}, which depend on the phase space variables of the system \cite{https://doi.org/10.1002/cpa.3160210503}. The dynamics of equation (\ref{eq:Hamilton}) is encoded in the equation

\begin{equation} \label{eq:Ham_Lax_Pairs}
    \dv{L}{t} \equiv \dot{L} = [M, L]
\end{equation}
where $[M, L]$ denotes the commutator of $M$ and $L$. 

The immediate interest in the existence of a Lax pair lies in its ability to generate conserved quantities in a straightforward way. Indeed, the solution to equation (\ref{eq:Ham_Lax_Pairs}) takes the form

\begin{equation}
    L(t) = g(t) L(0) g(t)^{-1}
\end{equation}
where the invertible matrix $g(t)$ satisfies
\begin{equation}
    M = \dv{g}{t} g^{-1}
\end{equation}
It follows that if $I(L)$ is a function invariant under conjugation, i.e., $L \mapsto gLg^{-1}$, then $I(L(t))$ is a constant of motion.
Such functions depend only on the eigenvalues of $L$. 
Therefore, the evolution equation (\ref{eq:Ham_Lax_Pairs}) is said to be \emph{isospectral}, meaning that the spectrum of $L$ is preserved by the time evolution.

\begin{remark}
    Recall that a system is integrable in the sense of Liouville if two conditions are met: the number of independent conserved quantities equals the number of degrees of freedom, and these conserved quantities are mutually in involution.
\end{remark}
\begin{remark}
    A Lax pair is not unique. In fact, even the size of the matrices can be modified. Moreover, there exists a natural gauge group acting on the Lax pair, defined by the transformation
    \begin{equation}
        L \longrightarrow gLg^{-1}, \quad M \longrightarrow gMg^{-1} + \dv{g}{t}g^{-1}
    \end{equation}
    where $g$ is an invertible matrix depending on the phase space variables.  
    This reflects the non-uniqueness of the Lax representation and is a particular case of the general gauge freedom of Lax pairs.
\end{remark}

\subsection{Example: Lax pairs for oscillators} \label{subsec:Example: Lax pairs for oscillators}
In this subsection we give a direct derivation of a Lax representation for one-dimensional systems starting from a simple ansatz.

We consider the following example.
\begin{enumerate} [label=\roman*)]
    \item Find a Lax pair representation $\dot{L} = [M, L]$ for the one-dimensional harmonic oscillator with Hamiltonian
    \begin{equation}
        H = \frac{p^2}{2} + \frac{\omega^2 q^2}{2}.
    \end{equation}
    Use the following ansatz for the $L$-operator:
    \begin{equation}
        L = \begin{bmatrix}
        p & f(q) \\
        f(q) & -p
        \end{bmatrix}.
    \end{equation}
    Determine the function $f(q)$ and the corresponding $M$ such that the Lax equation holds. 
    \item  How does the result change if the system is modified to an anharmonic oscillator instead of the harmonic one?
\end{enumerate}
\textbf{Solution i).}
From Hamilton's equations,
\begin{equation}
    \dot q = \pdv{H}{p} = p, \qquad
    \dot p = - \pdv{H}{q} = -\omega^2 q,
\end{equation}
we find that
\begin{equation}
    \dot L =
    \begin{bmatrix}
    \dot p & f'(q)\,\dot q \\
    f'(q)\, \dot q & -\dot p
    \end{bmatrix}
    =
    \begin{bmatrix}
    -\omega^2 q & f'(q)\,p \\
    f'(q)\,p & \omega^2 q
    \end{bmatrix}.
\end{equation}
We seek a Lax representation $\dot L = [M,L]$.
Let 
\begin{equation}
    M =
    \begin{bmatrix}
        a & b \\
        c & d
    \end{bmatrix}.
\end{equation}
Therefore
\begin{align}
    [M,L] = ML - LM &=
    \begin{bmatrix}
        ap + b f(q) & a f(q) - bp \\
        cp + d f(q) & c f(q) - dp
    \end{bmatrix} 
    - 
    \begin{bmatrix}
        pa + c f(q) & pb + d f(q) \\
        a f(q) - pc & b f(q) - pd
    \end{bmatrix} \\
     &=
    \begin{bmatrix}
        (b - c) f(q) & (a - d) f(q) - 2pb \\
        (d - a) f(q) + 2pc & (c - b) f(q)
    \end{bmatrix}.
\end{align}
Equating the two expressions, we obtain
\begin{equation}
    \dot L =
    \begin{bmatrix}
       -\omega^2 q & f'(q)\,p \\
        f'(q)\,p & \omega^2 q
    \end{bmatrix}
    = 
    \begin{bmatrix}
        (b - c) f(q) & (a - d) f(q) - 2pb \\
        (d - a) f(q) + 2pc & (c - b) f(q)
    \end{bmatrix}.
\end{equation}
Comparing the matrix elements we obtain
\begin{equation}\label{eq:system_lax_pari_ex}
    \begin{cases}
        - \omega^2 q = (b - c) f(q) \\
        f'(q) \, p = (a - d) f(q) - 2pb \\
        f'(q) \, p = (d - a) f(q) + 2pc
    \end{cases}.
\end{equation}
Since the equality must hold for all $p$, we equate separately the coefficients of $p$ and the terms independent of $p$:
\begin{equation}
    b = - c, \qquad d = a.
\end{equation}
From the relations obtained above, we have
\begin{equation}
    M = 
    \begin{bmatrix}
        a & b \\
        -b & a
    \end{bmatrix}
        = 
        a I + 
    \begin{bmatrix}
        0 & b \\
        -b & 0
    \end{bmatrix},
\end{equation}
where $ I $ denotes the identity matrix.
Since the identity commutes with any matrix, we have
\begin{equation}
    [L, aI] = LaI - aIL = a(LI - IL) = 0.
\end{equation}
Therefore, the term proportional to the identity does not contribute to the commutator $[M,L]$, and the Lax equation $\dot L = [M,L]$ remains unchanged under the transformation
\begin{equation}
    M \longrightarrow M - a I.
\end{equation}
We can thus choose, without loss of generality, $a = d = 0$ so that $M$ takes an antisymmetric form:
\begin{equation}
    M = 
    \begin{bmatrix}
        0 & b \\
        -b & 0
    \end{bmatrix}.
\end{equation}
This reflects the fact that $M$ is defined up to the addition of a multiple of the identity, which does not affect the commutator.
The system of equations (\ref{eq:system_lax_pari_ex}) becomes
\begin{equation}
    \begin{cases}
        - \omega^2 q = 2b f(q) \\
        f'(q)\, p = -2pb
    \end{cases}
    \implies
    \begin{cases}
        f(q) = -\frac{\omega^2}{2b} q \\
        f(q) = -2bq + K
    \end{cases}.
\end{equation}
Comparing the two expressions for $f(q)$, we obtain
\begin{equation}
    - \frac{\omega^2}{2b}q = - 2bq + K.
\end{equation}
Since this equality must hold for all values of $q$, we equate separately the coefficient of $q$ and the constant term. This gives
\begin{equation}
    \frac{\omega^2}{2b} = 2b,
    \qquad
    K = 0.
\end{equation}
This means that $b = \pm \omega/2$ and $f(q) = \pm \omega q$. The two choices correspond to equivalent Lax representations, so we fix the sign and take
\begin{equation}
    f(q) = \omega q, \qquad b = - \frac{\omega}{2}.
\end{equation}
Thus
\begin{equation}
    L = 
    \begin{bmatrix}
        p & \omega q \\
        \omega q & -p
    \end{bmatrix}
    \qquad
    M = 
    \begin{bmatrix}
        0 & -\omega / 2 \\
        \omega / 2 & 0
    \end{bmatrix}.
\end{equation}
\textbf{Solution ii).}
Let us now consider a more general one-dimensional system with Hamiltonian
\begin{equation}
    H(p,q) = \frac{p^2}{2} + V(q),
\end{equation}
whose equations of motion are
\begin{equation}
    \dot q = p, \qquad \dot p = -V'(q).
\end{equation}
We keep the same ansatz for the Lax pair,
\begin{equation}
    L =
    \begin{bmatrix}
        p & f(q) \\
        f(q) & -p
    \end{bmatrix},
    \qquad
    M =
    \begin{bmatrix}
        0 & b(q) \\
        -b(q) & 0
    \end{bmatrix},
\end{equation}
and require that the Lax equation $\dot L = [M,L]$
be satisfied.
The commutator has the same algebraic structure as before,
\begin{equation}
    [M,L] =
    \begin{bmatrix}
        2 b f & -2 p b \\
        -2 p b & -2 b f
    \end{bmatrix}.
\end{equation}
Using the equations of motion, the time derivative of $L$ is
\begin{equation}
    \dot L =
    \begin{bmatrix}
    \dot p & f'(q)\, \dot q \\
    f'(q) \, \dot q & -\dot p
    \end{bmatrix}
    =
    \begin{bmatrix}
    - V'(q) & f'(q) \, p \\
    f'(q) \, p & V'(q)
    \end{bmatrix}.
\end{equation}
Equating $\dot L = [M,L]$ componentwise gives
\begin{equation}
    \begin{cases}
        V'(q) = - 2 b f(q) \\
        f'(q) p = - 2 p b
    \end{cases}.
\end{equation}
From the second equation, valid for all $p$, we deduce
\begin{equation}
    b = - \frac{f'(q)}{2}.
\end{equation}
Substituting this into the first one yields
\begin{equation}
    V'(q) = f(q) f'(q).
\end{equation}
This can be rewritten as
\begin{equation}
    \frac{1}{2}\dv{}{q} f(q)^2 = V'(q),
\end{equation}
which integrates to
\begin{equation}
    f^2(q) = 2 V(q) + K,
\end{equation}
where $K$ is an integration constant. Whenever $2V(q)+K \ge 0$, this yields
\begin{equation}
    f(q) = \pm \sqrt{2 V(q) + K},
    \qquad
    b(q) = -\frac{V'(q)}{2f(q)}.
\end{equation}
For the harmonic oscillator $V(q) = \tfrac{1}{2}\omega^2 q^2$,
this reduces to $f(q) = \omega q$ and $b = - \omega/2$,
reproducing the previous result.

\section{The \texorpdfstring{$\boldsymbol{r}$}{r}-matrix structure}\label{sec:rmatrix_formalism}
In this section we present a self-contained derivation of the $r$-matrix structure and its main consequences, emphasizing explicit computations. The general framework follows the standard $r$-matrix formalism developed in \cite{SemenovTyanShanskii1983,Sklyanin1982,babelon2003introduction}.

A Lax pair provides a systematic way to construct conserved quantities, without making explicit use of a Poisson structure. However, Liouville integrability requires not only the existence of a Poisson structure, but also that the conserved quantities be in involution.
We now introduce a general form of the Poisson brackets between the entries of the Lax matrix which ensures this involution property.

Consider a Lax pair $L, M$ where both are $N \times N$ matrices, and assume that the matrix $L$ is diagonalizable, so that
\begin{equation} \label{eq:lambda_matrix}
    L = U \Lambda U^{-1}.
\end{equation}
The eigenvalues $\lambda_k$ of $L$ provide conserved quantities. We will not address here the question of their independence.

We introduce the following notation. Denote by $E_{ij}$ the canonical basis of the $ N \times N $ matrices, with components $(E_{ij})_{kl} = \delta_{ik} \delta_{jl}$. Then we can write
\begin{equation}
    L = \sum_{ij} L_{ij} E_{ij}.
\end{equation}
The coefficients $ L_{ij} $ of the Lax matrix are functions on the phase space. We can compute the Poisson brackets $ \{L_{ij}, L_{kl}\} $ and arrange them as follows. Define
\begin{equation}
    L_1 \equiv L \otimes 1 = \sum_{ij} L_{ij} (E_{ij} \otimes 1), 
    \quad
    L_2 \equiv 1 \otimes L = \sum_{ij} L_{ij} (1 \otimes E_{ij}).
\end{equation}
The subscript 1 or 2 indicates that the matrix $ L $ acts on the first or second factor of the tensor product, respectively. 
Likewise, for a tensor $ T $ belonging to the tensor product of two copies of $ N \times N $ matrices, we set
\begin{equation}
    T = T_{12} = \sum_{ijkl} T_{ijkl} E_{ij} \otimes E_{kl},
    \quad
    T_{21} = \sum_{ijkl} T_{ijkl} E_{kl} \otimes E_{ij}.
\end{equation}
More generally, when tensor products with several copies of $ N \times N $ matrices are involved, we denote by $ L_\alpha $ the embedding of $ L $ in the $\alpha$-th position, e.g.
\begin{equation}
    L_3 = 1 \otimes 1 \otimes L \otimes \cdots,
    \quad
    T_{\alpha\beta} \text{ the embedding of } T \text{ in positions } \alpha \text{ and } \beta.
\end{equation}
We also use $ \Tr_\alpha $ to denote the partial trace over the $\alpha$-th factor in a tensor product. For instance,
\begin{equation}
    \Tr_1 T_{12} = \sum_{ijkl} T_{ijkl} \, \Tr(E_{ij}) \, E_{kl}.
\end{equation}
\begin{definition}
    Define $ \{L_1, L_2\} $ as the matrix of Poisson brackets between the entries of $ L $:
    \begin{equation}
        \{L_1, L_2\} = \sum_{ijkl} \{L_{ij}, L_{kl}\} E_{ij} \otimes E_{kl}.
    \end{equation}
\end{definition}
For integrable systems, the Poisson brackets between the entries of the Lax matrix admit a particularly simple and structured form.

We collect a number of auxiliary results that will be used in the formulation of the $r$-matrix structure. We start with a general identity, which will serve as the basis for the subsequent lemmas and corollaries.

\begin{proposition}\label{prop:leibniz-expand-general}
    Let $L$ and $X$ be matrices whose entries belong to a Poisson algebra.
    Set $L_1 = L \otimes 1$ and $X_2 = 1 \otimes X$.
    Assume that the Poisson bracket is extended entrywise to matrix-valued functions and satisfies the Leibniz rule in each argument.
    Then, for every integer $n \ge 1$,
    \begin{equation}\label{eq:leibniz-expand-general}
        \{L_1^n, X_2\}
        =
        \sum_{p=0}^{n-1}
        L_1^{n-1-p} \{L_1, X_2\} L_1^p.
    \end{equation}
\end{proposition}

\begin{proof}
    We argue by induction on $n$.
    For $n=1$, the statement is immediate.
    Assume that \eqref{eq:leibniz-expand-general} holds for some $n$. Then, by the Leibniz rule,
    \begin{equation}
        \{L_1^{n+1}, X_2\} = \{L_1 L_1^n, X_2\} = L_1 \{L_1^n, X_2\} + \{L_1, X_2\} L_1^n.
    \end{equation}
    Using the inductive hypothesis, we obtain
    \begin{equation}
        \{L_1^{n+1}, X_2\}
        = \sum_{p=0}^{n-1} L_1^{n-p} \{L_1, X_2\} L_1^p + \{L_1, X_2\} L_1^n.
    \end{equation}
    The last term is precisely the contribution corresponding to $p=n$, hence
    \begin{equation}
        \{L_1^{n+1}, X_2\}
        = \sum_{p=0}^{n} L_1^{n-p} \{L_1, X_2\} L_1^p.
    \end{equation}
    This proves the inductive step and therefore the statement for all $n \ge 1$.
\end{proof}

\begin{corollary}\label{cor:leibniz-expand}
    Let $L$ be a matrix whose entries belong to a Poisson algebra.
    Set $L_1 = L \otimes 1$ and $L_2 = 1 \otimes L$.
    Assume that the Poisson bracket is extended entrywise to matrix-valued functions and satisfies the Leibniz rule in each argument.
    Then, for all integers $n,m \ge 1$,
    \begin{equation}\label{eq:leibniz-expand}
        \{L_1^n,L_2^m\}
        =
        \sum_{p=0}^{n-1}\sum_{q=0}^{m-1}
        L_1^{\,n-1-p}L_2^{\,m-1-q}
        \{L_1,L_2\}
        L_1^pL_2^q
    \end{equation}
\end{corollary}

\begin{proof}
    Applying Proposition \ref{prop:leibniz-expand-general} with $X_2 = L_2^m$, we obtain
    \begin{equation}
        \{L_1^n, L_2^m\}
        =
        \sum_{p=0}^{n-1}
        L_1^{n-1-p} \{L_1, L_2^m\} L_1^p.
    \end{equation}
    Expanding $\{L_1, L_2^m\}$ by the Leibniz rule in the second argument gives
    \begin{equation}
        \{L_1, L_2^m\}
        =
        \sum_{q=0}^{m-1}
        L_2^{m-1-q} \{L_1, L_2\} L_2^q.
    \end{equation}
    Substituting this into the previous formula yields
    \begin{equation}
        \{L_1^n, L_2^m\}
        =
        \sum_{p=0}^{n-1} \sum_{q=0}^{m-1}
        L_1^{n-1-p} L_2^{m-1-q}
        \, \{L_1, L_2\} \,
        L_1^p L_2^q,
    \end{equation}
    which proves the claim.
\end{proof}

\begin{proposition}\label{prop:fund}
     Let $L$ be a matrix whose entries belong to a Poisson algebra.
    Set $L_1 = L \otimes 1$ and $L_2 = 1 \otimes L$.
    Assume that the Poisson bracket is extended entrywise to matrix-valued functions and satisfies the Leibniz rule in each argument.
    Suppose that there exists an element $r$ of the tensor product algebra such that
    \begin{equation}\label{eq:fund}
        \{L_1,L_2\} = [r_{12},L_1] - [r_{21},L_2].
    \end{equation}
    Then, for all integers $n, m \ge 1$,
    \begin{equation}\label{eq:claim}
        \{L_1^n, L_2^m\} = [a_{12}^{n,m}, L_1] + [b_{12}^{n,m}, L_2],
    \end{equation}
    where
    \begin{align}
        a_{12}^{n,m} &=\sum_{p=0}^{n-1}\sum_{q=0}^{m-1} L_1^{n-1-p}L_2^{m-1-q}r_{12}L_1^{p}L_2^{q}, \\
        b_{12}^{n,m} &=-\sum_{p=0}^{n-1}\sum_{q=0}^{m-1} L_1^{n-1-p}L_2^{m-1-q}r_{21}L_1^{p}L_2^{q}.
    \end{align}
\end{proposition}
\begin{proof}
   We start from the result of the previous corollary~\ref{cor:leibniz-expand}:
    \begin{equation}
        \{ L_1^n, L_2^m \} = \sum_{p=0}^{n-1} \sum_{q=0}^{m-1} L_1^{n-1-p} L_2^{m-1-q} \{ L_1, L_2 \} L_1^{p} L_2^{q}.
    \end{equation}
    Using Corollary~\ref{cor:leibniz-expand} and substituting \eqref{eq:fund}, we obtain
    \begin{equation}
        \{ L_1^n, L_2^m \}
        =\sum_{p=0}^{n-1} \sum_{q=0}^{m-1}
        L_1^{n-1-p} L_2^{m-1-q} \left( [r_{12}, L_1] - [r_{21}, L_2] \right) L_1^{p} L_2^{q}.
    \end{equation}
    Split the sum into the two corresponding contributions. For the $r_{12}$-term we have
    \begin{align}
        \sum_{p,q} L_1^{n-1-p} L_2^{m-1-q} [r_{12}, L_1] L_1^{p} L_2^{q} &= \sum_{p,q} \left( L_1^{n-1-p} L_2^{m-1-q} r_{12} L_1^{p+1} L_2^{q} - L_1^{n-p} L_2^{m-1-q} r_{12} L_1^{p} L_2^{q} \right)\\
        &= \left[ \sum_{p,q} L_1^{n-1-p}L_2^{m-1-q} r_{12} L_1^{p}L_2^{q}\;, \;L_1 \right],
    \end{align}
    where in the second line we have used that $L_1$ and $L_2$ commute, so that the factors $L_1^{n-1-p}L_2^{m-1-q}$ and $L_1^{p}L_2^{q}$ commute with $L_1$. Hence
    \begin{equation}
        X [r_{12},L_1] Y = [X r_{12} Y, L_1],
    \end{equation}
    and we obtain
    \begin{equation}
        a_{12}^{n,m} = \sum_{p=0}^{n-1}\sum_{q=0}^{m-1}
        L_1^{n-1-p} L_2^{m-1-q} r_{12} L_1^{p} L_2^{q}.
    \end{equation}
    An analogous computation for the $-[r_{21},L_2]$ term gives
    \begin{equation}
    \sum_{p,q} L_1^{n-1-p} L_2^{m-1-q} (-[r_{21},L_2]) L_1^{p} L_2^{q}
    =
    \left[
    -\sum_{p,q} L_1^{n-1-p} L_2^{m-1-q} r_{21} L_1^{p} L_2^{q},
    L_2
    \right].
\end{equation}
Combining the two commutators yields \eqref{eq:claim}.


    
\end{proof}

\begin{lemma} \label{lemma:TrL^n involution}
    If (\ref{eq:claim}) holds then the quantities $\Tr(L^n)$ are in involution, i.e.
    \begin{equation}
        \{\Tr(L^n), \Tr(L^m)\} = 0 \quad \forall \, n, m.
    \end{equation}
\end{lemma}
\begin{proof}
    Take the trace over both spaces:
    \begin{equation}
        \{\Tr\left(L^n\right) , \Tr\left(L^m\right)\}
        = \Tr_{12}\left(\left\{L_1^n , L_2^m\right\}\right).
    \end{equation}
    Using the linearity 
    \begin{equation}
        \{\Tr\left(L^n\right) , \Tr\left(L^m\right)\} 
        = \Tr_{12}\left(\left[a_{12}^{n,m} , L_1\right]\right) + \Tr_{12}\left(\left[b_{12}^{n,m} , L_2\right]\right).
    \end{equation}
    By cyclicity of the trace
    \begin{equation}
        \Tr_{12}\left(\left[X , Y\right]\right) = 0,
    \end{equation}
    hence
    \begin{equation}
        \{\Tr\left(L^n\right) , \Tr\left(L^m\right)\} = 0,
    \end{equation}
    this shows that the spectral invariants $\Tr(L^n)$ are in involution.
\end{proof}


\begin{corollary} \label{cor:importante}
    Assume there exists an element $r_{12}$ of the tensor product algebra such that
    \begin{equation}
        \{L_1, L_2\} = [r_{12}, L_1] - [r_{21}, L_2].
    \end{equation}
    Then the spectral invariants $\Tr (L^n)$ are in involution, i.e.
    \begin{equation}
        \{\Tr(L^n), \Tr(L^m)\} = 0 \quad \, \forall \, n,m.
    \end{equation}
\end{corollary}

\begin{proof}
    This follows directly from Proposition~\ref{prop:fund} together with Lemma~\ref{lemma:TrL^n involution}.
\end{proof}

We now address the converse problem: starting from the involution of the eigenvalues of $L$, we show that, locally on phase space, one can reconstruct an $r$-matrix structure of the form (\ref{eq:r_matrix}).

\begin{proposition}[Local reconstruction of an $r$-matrix]\label{prop:existence_rmatrix}
    Assume that $L$ is locally diagonalizable as in (\ref{eq:lambda_matrix}) (for instance, with simple spectrum), and that its eigenvalues Poisson commute,
    \begin{equation}
        \{\lambda_i,\lambda_j\}=0 \quad \forall \, i,j.
    \end{equation}
    Then, locally on phase space, one can reconstruct an element $r_{12}$ of the tensor product algebra such that
    \begin{equation}\label{eq:r_matrix}
        \{L_1, L_2\} = [r_{12}, L_1] - [r_{21}, L_2].
    \end{equation}
\end{proposition}

\begin{proof}[Sketch of proof]
We outline the construction locally on a region of phase space where $L$ is diagonalizable with simple spectrum, so that one can choose $U$ smoothly such that $L = U \Lambda U^{-1}$.

Using the Leibniz rule, we expand
\begin{equation}
    \{L_1,L_2\} = \{U_1 \Lambda_1 U_1^{-1}, U_2 \Lambda_2 U_2^{-1}\}.
\end{equation}
All terms can be expressed in terms of $U$, $\Lambda$, and the brackets $\{U_1,U_2\}$, $\{U_1,\Lambda_2\}$, $\{\Lambda_1,U_2\}$, using the identity
\begin{equation}
    \{U^{-1},X\} = - U^{-1} \{U,X\} U^{-1}.
\end{equation}

Conjugating the whole expression by $U_1^{-1} \otimes U_2^{-1}$, one reduces to computing the Poisson brackets of $\Lambda$ together with the transformed brackets of $U$. The terms involving $\{\Lambda_1,\Lambda_2\}$ vanish by assumption, since the eigenvalues Poisson commute.

The remaining contributions can be organized into commutators with $\Lambda_1$ and $\Lambda_2$, yielding an expression of the form
\begin{equation}
    \{L_1,L_2\} = [r_{12},L_1] - [r_{21},L_2],
\end{equation}
where $r_{12}$ is constructed locally from $U$ and its Poisson brackets.

The construction is therefore valid on any open set where the diagonalization of $L$ can be chosen smoothly.
\end{proof}


\section{Commuting flows}\label{sec:Commuting flows}
The Poisson brackets~\eqref{eq:r_matrix} imply that the spectral invariants of $L$ are in involution \cite{babelon2003introduction}. A convenient set of commuting Hamiltonians is given by
\begin{equation}
    H_n = \Tr\left(L^n\right).
\end{equation}
The Hamiltonians $H_n$ are in involution, $\{H_n,H_m\}=0$, as a consequence of the $r$-matrix structure. Furthermore, we show that the time evolution of the Lax matrix $L$ with Hamiltonian $H_n$ naturally takes the Lax form.
\begin{proposition}
    Suppose that $\{L_1,L_2\} = [r_{12},L_1] - [r_{21},L_2]$.
    If we take $H_n = \Tr\left(L^n\right)$ as Hamiltonians, then the
    equations of motion admit a Lax representation:
    \begin{equation}
        \dv{L}{t_n} = \{L,H_n\} = [M_n,L], \quad
        \text{ with } \quad M_n = -n \, \Tr_1 \left( L_1^{n-1} r_{21} \right).
    \end{equation}
\end{proposition}
\begin{proof}
    We set $m = 1$ in equation~\eqref{eq:claim} obtaining
    \begin{equation}
        \{L_1^n,L_2\} = [a_{12}^{n,1},L_1] + [b_{12}^{n,1},L_2],
    \end{equation}
    where
    \begin{equation}
        b_{12}^{n,1} = -\sum_{p=0}^{n-1} L_1^{\,n-1-p} r_{21} L_1^p.
    \end{equation}
    Taking the trace over the first tensor factor, the first commutator vanishes by cyclicity of the trace, hence
    \begin{equation}
        \{L_2,H_n\} = \Tr_1\left(\{L_1^n,L_2\}\right) = \left[\Tr_1\left(b_{12}^{n,1}\right),L_2\right].
    \end{equation}
    Identifying $L_2 = 1 \otimes L$, we obtain
    \begin{equation}
        \{L,H_n\} = \left[\Tr_1\left(b_{12}^{n,1}\right),L\right].
    \end{equation}
    By cyclicity of the trace in the first factor, each term contributes equally:
    \begin{equation}
        \Tr_1\left(L_1^{\,n-1-p} r_{21} L_1^p\right)
        = \Tr_1\left(L_1^{\,n-1} r_{21}\right).
    \end{equation}
    Therefore
    \begin{equation}
        \Tr_1\left(b_{12}^{n,1}\right)
        = -n\,\Tr_1\left(L_1^{\,n-1} r_{21}\right),
    \end{equation}
    and thus
    \begin{equation}
        \{H_n,L\} = [M_n,L],
    \end{equation}
    with
    \begin{equation}
        M_n = -n\,\Tr_1\left(L_1^{\,n-1} r_{21}\right).
    \end{equation}
\end{proof}

\begin{remark}
    Note that the matrices $M_n$ are not unique since adding any matrix commuting with $L$ does not change the equations of motion.
\end{remark}


\section{Tensor Casimir element properties}\label{sec:casimir}
In this section we develop the properties of the tensor Casimir element that are needed in the $r$-matrix formalism \cite{perelomov1990integrable,babelon2003introduction}.

The $r$-matrix formalism is naturally expressed in terms of tensorial objects associated with Lie algebras. Among these, a distinguished role is played by the Casimir element. It provides a canonical invariant tensor and will be used to constrain the structure of $r$-matrices.

We now recall its definition and main properties.

\begin{definition}
    The Casimir tensor element of $\mathfrak{gl}(n)$ associated with the trace form is defined as
    \begin{equation}\label{eq:casimir_element}
        C_{12} = \sum_{i,j} E_{ij} \otimes E_{ji}.
    \end{equation}
\end{definition}
\begin{proposition}\label{prop:casimir_invariant}
    The Casimir tensor element satisfies
    \begin{equation}\label{eq:casimir_property}
        [C_{12}, X \otimes 1 + 1 \otimes X] = 0
        \qquad \forall X \in \mathfrak{gl}(n).
    \end{equation}
\end{proposition}

\begin{proof}
Let $X_1=X\otimes1$ and $X_2=1\otimes X$, and write
\begin{equation}
    X=\sum_{k,l}X_{kl}E_{kl}.
\end{equation}
By equation (\ref{eq:casimir_element}) we have
\begin{equation}
    [C_{12},X_1]
    =\sum_{i,j}[E_{ij},X]\otimes E_{ji} \quad \text{ and } \quad [C_{12},X_2]
    =\sum_{i,j}E_{ij}\otimes[E_{ji},X].
\end{equation}
Using
\begin{equation}\label{eq:[E_ij,E_kl]}
    [E_{ij},E_{kl}]
    =\delta_{jk}E_{il}-\delta_{il}E_{kj},
\end{equation}
we obtain
\begin{equation}
    [E_{ij},X]
    =\sum_l X_{jl}E_{il}-\sum_k X_{ki}E_{kj} \quad \text{ and } \quad 
    [E_{ji},X]
    =\sum_l X_{il}E_{jl}-\sum_k X_{kj}E_{ki}.
\end{equation}
Thus
\begin{align}
[C_{12},X_1]
    &=\sum_{i,j,l}X_{jl}E_{il}\otimes E_{ji}
    -\sum_{i,j,k}X_{ki}E_{kj}\otimes E_{ji}, \\
    [C_{12},X_2]
    &=\sum_{i,j,l}X_{il}E_{ij}\otimes E_{jl}
    -\sum_{i,j,k}X_{kj}E_{ij}\otimes E_{ki}.
\end{align}
After relabelling dummy indices, the first term equals the fourth and the second equals the third with opposite sign. Therefore
\begin{equation}
    [C_{12},X_1]+[C_{12},X_2]=0.
\end{equation}
\end{proof}

\begin{lemma}[Flip property of the Casimir tensor element]\label{lemma:flipping}
    Let
    \begin{equation}
        C_{12} = \sum_{i,j} E_{ij} \otimes E_{ji}
    \end{equation}
    be the Casimir tensor element of $\mathfrak{gl}(n)$.
    Then, for all $X,Y \in \mathfrak{gl}(n)$, one has
    \begin{equation}
        C_{12}(X \otimes Y)C_{12} = Y \otimes X.
    \end{equation}
\end{lemma}

\begin{proof}
    Write
    \begin{equation}
        X = \sum_{a,b} X_{ab} E_{ab},
        \qquad
        Y = \sum_{c,d} Y_{cd} E_{cd}.
    \end{equation}
    Then
    \begin{equation}
        C_{12}(X \otimes Y)C_{12}
        =
        \sum_{i,j,k,l}
        (E_{ij} X E_{kl}) \otimes (E_{ji} Y E_{lk}).
    \end{equation}
    Using
    \begin{equation}
        E_{ij} E_{ab} E_{kl}
        =
        \delta_{ja}\delta_{bk} E_{il},
        \qquad
        E_{ji} E_{cd} E_{lk}
        =
        \delta_{ic}\delta_{dl} E_{jk},
    \end{equation}
    we obtain
    \begin{equation}
        E_{ij} X E_{kl}
        =
        \sum_{a,b} X_{ab}\delta_{ja}\delta_{bk} E_{il}
        =
        X_{jk} E_{il},
    \end{equation}
    and similarly
    \begin{equation}
        E_{ji} Y E_{lk}
        =
        \sum_{c,d} Y_{cd}\delta_{ic}\delta_{dl} E_{jk}
        =
        Y_{il} E_{jk}.
    \end{equation}
    Therefore
    \begin{equation}
        C_{12}(X \otimes Y)C_{12}
        =
        \sum_{i,j,k,l}
        X_{jk} Y_{il}\,
        E_{il} \otimes E_{jk}.
    \end{equation}
    Rearranging the sums, we get
    \begin{equation}
        C_{12}(X \otimes Y)C_{12}
        =
        \left(\sum_{i,l} Y_{il} E_{il}\right)
        \otimes
        \left(\sum_{j,k} X_{jk} E_{jk}\right)
        =
        Y \otimes X.
    \end{equation}
\end{proof}

The following identity will play a crucial role in the derivation of the Yang--Baxter equation in Section~\ref{sec:YB}.
\begin{proposition}\label{prop:casimir_r_relation}
    Let $C_{12}$ be the Casimir tensor element and let $r \in \mathfrak{gl}(n)^{\otimes 2}$. Then
    \begin{equation}\label{eq:casimir_YB}
        [C_{12},r_{13}] + [C_{12},r_{23}]
        + [C_{13},r_{21}] + [C_{13},r_{23}]
        + [C_{23},r_{21}] + [C_{23},r_{31}] = 0.
    \end{equation}
\end{proposition}
\begin{proof}
    We rewrite the left-hand side of \eqref{eq:casimir_YB} as
    \begin{equation}\label{eq:casimir_accorpati}
        [C_{12}, r_{13} + r_{23}]
        + [C_{13}, r_{21} + r_{23}]
        + [C_{23}, r_{21} + r_{31}].
    \end{equation}
    We consider the first term. By Lemma~\ref{lemma:flipping}, the operator $C_{12}$ acts as the transposition of the first two tensor factors. Therefore
    \begin{equation}
        C_{12} r_{13} C_{12} = r_{23},
        \qquad
        C_{12} r_{23} C_{12} = r_{13}.
    \end{equation}
    Hence
    \begin{equation}
        C_{12}(r_{13}+r_{23})C_{12} = r_{13}+r_{23}.
    \end{equation}
    Since $C_{12}^2 = 1$, then
    \begin{equation}
        C_{12}^{-1} (r_{13} + r_{23}) C_{12} = r_{13} + r_{23} \implies [C_{12}, r_{13}+r_{23}] = 0.
    \end{equation}
    The same argument applied cyclically gives
    \begin{equation}
        [C_{13}, r_{21}+r_{23}] = 0,
        \qquad
        [C_{23}, r_{21}+r_{31}] = 0.
    \end{equation}
    Summing these three identities, we obtain \eqref{eq:casimir_YB}.
\end{proof}

\section{Classical and Modified Yang--Baxter equation}\label{sec:YB}
We now give a direct treatment of the classical and modified Yang–Baxter equations in the present framework \cite{SemenovTyanShanskii1983,Sklyanin1982}.

The $r$-matrix formalism can be formulated in a natural way in the setting of Lie algebras. In particular, the classical Yang--Baxter equation is most conveniently expressed in terms of tensor products of a Lie algebra and its invariant structures.

We therefore introduce the general algebraic framework that will be used throughout this section.

\begin{definition}
    Let $\mathcal{G}$ be a finite-dimensional Lie algebra over $\mathbb{C}$, equipped with an invariant non-degenerate bilinear form $(\cdot,\cdot)$. We denote by $\mathcal{G}^{\otimes n}$ the $n$-fold tensor product of $\mathcal{G}$, and use the standard notation $X_1, X_2, X_3$ for the embedding of $X \in \mathcal{G}$ into the corresponding tensor factors.
\end{definition}

\begin{definition}\label{def:double_bracket}
    Let $r \in \mathcal{G}^{\otimes 2}$. We define
    \begin{equation}
        [[r,r]] \equiv [r_{12},r_{13}] + [r_{12},r_{23}] + [r_{13},r_{23}]
        \in \mathcal{G}^{\otimes 3}.
    \end{equation}
\end{definition}
\begin{definition}
    Let $C \in \mathcal{G}^{\otimes 2}$ be the Casimir tensor element.  
    We define
    \begin{equation}
        \Omega \equiv [[C,C]] \in \mathcal{G}^{\otimes 3}.
    \end{equation}
\end{definition}
\begin{remark}
    The tensor $\Omega$ is called the canonical invariant element in $\mathcal{G}^{\otimes 3}$.
\end{remark}
\begin{definition}
    Let $r \in \mathcal{G}^{\otimes 2}$. The \emph{classical Yang--Baxter equation} (CYBE) is
    \begin{equation} \label{eq:CYBE}
        [[r,r]] = 0,
    \end{equation}
    while the \emph{modified classical Yang--Baxter equation} (mCYBE) is
    \begin{equation} \label{eq:mCYBE}
        [[r,r]] = \mu \, \Omega,
    \end{equation}
    where $\mu \in \mathbb{C}$ is a constant.
\end{definition}

\begin{proposition}
    Let $L$ be a Lax matrix whose Poisson brackets satisfy \eqref{eq:fund}, and assume that $r \in \mathcal{G}^{\otimes 2}$ is an antisymmetric constant tensor. If $r$ satisfies the modified classical Yang--Baxter equation, then the Jacobi identity for the Poisson bracket is satisfied:
    \begin{equation}
        \{\{L_1,L_2\},L_3\}
        + \{\{L_2,L_3\},L_1\}
        + \{\{L_3,L_1\},L_2\}
        = 0.
    \end{equation}
\end{proposition}
\begin{proof}
    We directly compute the terms in the Jacobi identity:
    \begin{equation} \label{eq:L_1 L_2 L_3}
        \{ \{L_1 , L_2 \} , L_3 \} \overset{(\ref{eq:fund})}{=} \{ [r_{12}, L_1] - [r_{21}, L_2] , L_3 \} = \{ [r_{12}, L_1] , L_3 \} - \{ [r_{21}, L_2] , L_3 \}.
    \end{equation}
    Since $r$ is assumed to be constant, it Poisson-commutes with $L$, and since it is antisymmetric we have
    \begin{equation}
        r_{21} = -r_{12}, \qquad r_{32} = -r_{23}, \qquad r_{31} = -r_{13}.
    \end{equation}
    We focus on 
    \begin{equation}
        \{ [r_{12}, L_1] , L_3 \} = \{ r_{12} L_1 - L_1 r_{12}, L_3 \} = r_{12} \{ L_1, L_3 \} - \{ L_1, L_3 \} r_{12} = [r_{12}, \{ L_1, L_3 \}],
    \end{equation}
    but $\{ L_1, L_3 \} = [r_{13}, L_1 ] - [r_{31}, L_3 ]$ then 
    \begin{equation}
        \{ [r_{12}, L_1] , L_3 \} = [r_{12}, [r_{13}, L_1]] - [r_{12}, [r_{31}, L_3]],
    \end{equation}
    in the same way we can compute the other term
    \begin{equation}
        \{ [r_{21}, L_2] , L_3 \} = [r_{21}, [r_{23}, L_2]] - [r_{21}, [r_{32}, L_3]].
    \end{equation}
    Hence \eqref{eq:L_1 L_2 L_3} reduces to the sum of 
    \begin{align} 
        \{ \{L_1, L_2 \}, L_3 \} &= [r_{12}, [r_{13}, L_1]] - [r_{12}, [r_{31}, L_3]] - [r_{21}, [r_{23}, L_2]] + [r_{21}, [r_{32}, L_3]] \label{eq:L_123},\\
        \{ \{L_3, L_1 \}, L_2 \} &= [r_{31}, [r_{32}, L_3]] - [r_{31}, [r_{23}, L_2]] - [r_{13}, [r_{12}, L_1]] + [r_{13}, [r_{21}, L_2]] \label{eq:L_321},\\
        \{ \{L_2, L_3 \}, L_1 \} &= [r_{23}, [r_{21}, L_2]] - [r_{23}, [r_{12}, L_1]] - [r_{32}, [r_{31}, L_3]] + [r_{32}, [r_{13}, L_1]] \label{eq:L_231}.
    \end{align}
    Summing \eqref{eq:L_123}, \eqref{eq:L_321}, and \eqref{eq:L_231}, we collect separately the terms proportional to $L_1$, $L_2$, and $L_3$.
    The terms containing $L_1$ are
    \begin{equation}
        [r_{12},[r_{13},L_1]] - [r_{13},[r_{12},L_1]] - [r_{23},[r_{12},L_1]] + [r_{32},[r_{13},L_1]].
    \end{equation}
    Using the Jacobi identity in the form $[a,[b,X]] - [b,[a,X]] = [[a,b],X]$  the first two terms give
    \begin{equation}
        [[r_{12},r_{13}],L_1].
    \end{equation}
    While for the last two terms we use Jacobi again to obtain
    \begin{equation}
        - [r_{23},[r_{12},L_1]] =
        - [[r_{23},r_{12}],L_1] - [r_{12},[r_{23},L_1]],
    \end{equation}
    and
    \begin{equation}
        [r_{32},[r_{13},L_1]] =
        [[r_{32},r_{13}],L_1] + [r_{13},[r_{32},L_1]].
    \end{equation}
    Hence the full contribution proportional to $L_1$ becomes
    \begin{equation} \label{eq:full_L1}
     [[r_{12},r_{13}],L_1] + [[r_{12},r_{23}],L_1] + [[r_{13},r_{23}],L_1] - [r_{12},[r_{23},L_1]] + [r_{13},[r_{32},L_1]].
    \end{equation}
    Since $r_{23}$ and $r_{32}$ act on the second and third tensor factors, while $L_1$ acts on the first one, we have
    \begin{equation}
        [r_{23},L_1] = [r_{32},L_1] = 0 \implies - [r_{12},[r_{23},L_1]] = 0 = [r_{13},[r_{32},L_1]],
    \end{equation}
    hence the residual terms vanish.
    Thus the full contribution proportional to $L_1$ becomes
    \begin{equation}
        [[r_{12},r_{13}] + [r_{12},r_{23}] + [r_{13},r_{23}], L_1].
    \end{equation}
    Using the modified classical Yang--Baxter equation, this is equal to
    \begin{equation}
        \mu [\Omega, L_1].
    \end{equation}
    Summing the three contributions, we obtain
    \begin{equation}
        \mu \Big( [\Omega,L_1] + [\Omega,L_2] + [\Omega,L_3] \Big)
        = \mu [\Omega, L_1 + L_2 + L_3].
    \end{equation}
    Now, by definition,
    \begin{equation}
        \Omega = [C_{12},C_{13}] + [C_{12},C_{23}] + [C_{13},C_{23}].
    \end{equation}
    
    Each of the tensors $C_{12}$, $C_{13}$ and $C_{23}$ commutes with $L_1+L_2+L_3$. 
    Indeed, by invariance of the Casimir tensor element, $C_{12}$ commutes with $L_1+L_2$, and it trivially commutes with $L_3$; similarly, $C_{13}$ commutes with $L_1+L_3$, and $C_{23}$ commutes with $L_2+L_3$. Therefore
    \begin{equation}
        [C_{12},L_1+L_2+L_3]=[C_{13},L_1+L_2+L_3]=[C_{23},L_1+L_2+L_3]=0.
    \end{equation}
    It follows that each commutator appearing in $\Omega$ also commutes with $L_1+L_2+L_3$, hence
    \begin{equation}
        [\Omega, L_1+L_2+L_3] = 0.
    \end{equation}
\end{proof}

\begin{lemma}
    Let $r_{12} \in \mathcal{G}^{\otimes2}$ be an antisymmetric solution of the modified classical Yang--Baxter equation, i.e.
    \begin{equation}
        r_{21} = -r_{12}, \qquad [[r,r]] = \mu \, \Omega .
    \end{equation}
    Let $C_{12}$ be the Casimir tensor element. For a choice of square root of $-\mu$, define
    \begin{equation}
        r_{12}^{\pm} = r_{12} \pm \sqrt{-\mu}\, C_{12}.
    \end{equation}
    Then $r_{12}^{\pm}$ are solutions of the classical Yang--Baxter equation:
    \begin{equation}\label{eq:cyb_pm}
        [r_{12}^{\pm},r_{13}^{\pm}] + [r_{12}^{\pm},r_{23}^{\pm}] + [r_{13}^{\pm},r_{23}^{\pm}] = 0.
    \end{equation}
\end{lemma}
\begin{proof}
    Substituting $r_{12}^{\pm}=r_{12}\pm \sqrt{-\mu} \, C_{12}$ into \eqref{eq:cyb_pm} and expanding, we obtain
    \begin{align}
        [r_{12},r_{13}] + [r_{12},&r_{23}] + [r_{13},r_{23}] \notag \\
        \pm \sqrt{-\mu} \bigg( 
        [C_{12},r_{13}] + [C_{12},r_{23}]
        + [r_{12},C_{13}] +& [r_{12},C_{23}]
        + [C_{13},r_{23}] + [r_{13},C_{23}]
        \bigg) \notag \\
        -\mu \bigg(
        [C_{12},C_{13}] + [C_{12},&C_{23}] + [C_{13},C_{23}]
        \bigg)
    \end{align}
    By \eqref{eq:casimir_YB}, the second term vanishes. Moreover, since $r_{12}$ satisfies the modified classical Yang--Baxter equation \eqref{eq:mCYBE}, the statement follows.
\end{proof}


\section{The open Toda chain}
The algebraic structures introduced in Sections~\ref{sec:casimir} and~\ref{sec:YB}
provide the general framework for the $r$-matrix formalism.
We now show how these structures are realized concretely in the Toda chain, which gives a fundamental example of an integrable system associated with a simple Lie algebra \cite{toda1967vibration,Olshanetsky1979,KOSTANT1979195,perelomov1990integrable}.
\subsection{The model}\label{subsec:The model}
The open Toda chains are integrable systems associated with Lie algebras.
Let $\mathcal{G}$ be a finite-dimensional simple Lie algebra of rank $r$, equipped with an invariant scalar product denoted by $( \cdot, \cdot )$. We choose a Cartan subalgebra with an orthonormal basis $\{H_i\}$. We have the Cartan decomposition of $\mathcal{G}$:
\begin{equation}
    \mathcal{G} = \mathcal{N}_{-} \oplus \mathcal{H} \oplus \mathcal{N}_{+} \quad \text{with} \quad \mathcal{N}_{\pm} = \bigoplus_{\beta \text{ positive}} \mathcal{G}_{\pm \beta},
\end{equation}
where $\mathcal{G}_{\pm \beta}$ are one-dimensional, generated by the root vectors $E_{\pm \beta}$ for any positive root $\beta$.
We will need the commutation relations:
\begin{align}
    [H_i, H_j] &= 0, \\
    [H, E_{\pm \alpha}] &= \pm \alpha(H) \, E_{\pm \alpha}, \\
    [E_{\alpha}, E_{-\alpha}] &= ( E_{\alpha}, E_{-\alpha} ) H_\alpha,
\end{align}
where $H_{\alpha} = \sum_i \alpha(H_i) \, H_i$. We will often use the constants $n_\alpha$ such that 
\begin{equation}
    n_\alpha^2 ( E_{\alpha},E_{-\alpha} ) = 1.
\end{equation}
Toda chains are associated with any simple Lie algebra $\mathcal{G}$ as follows. Consider two elements in the Cartan subalgebra:
\begin{equation}
    q = \sum_{i = 1}^r q_i H_i, \quad p = \sum_{i = 1}^r p_i H_i \quad r = \rank \mathcal{G}.
\end{equation}
The coefficients $(q_i, p_i)$ are the coordinates of a $2r$-dimensional phase space with Poisson brackets 
\begin{equation}\label{eq:poison_bis}
    \{q_i, p_j\} = \frac{1}{2} \delta_{ij},
\end{equation}
where the factor $1/2$ is introduced to simplify later formulae.
The Hamiltonian of the open Toda chain is by definition:
\begin{equation}
    H = ( p, p ) + 2 \sum_{\alpha \text{ simple}} \mathrm{e}^{2\alpha(q)}.
\end{equation}
The equations of motion are then:
\begin{align}
    \dv{q}{t} &= \{ q , H \} = p \\
    \dv{p}{t} &= \{ p , H \} = - 2 \sum_{\alpha \text{ simple}} H_\alpha \mathrm{e}^{2\alpha(q)}
\end{align}
or, combining the two,
\begin{equation}
    \dv[2]{q}{t} = - 2 \sum_{\alpha \text{ simple}} H_\alpha \mathrm{e}^{2\alpha(q)}.
\end{equation}

In the usual $\mathfrak{sl}(r+1)$ case, the Lie algebra of traceless $(r+1) \times (r+1)$ matrices, the Cartan algebra can be taken as traceless diagonal matrices.
The simple root vectors $E_{\alpha_i}$ are the canonical matrices $E_{i, i+1}$.
In this case it is not very convenient to use an orthonormal basis of the Cartan algebra under the scalar product $ ( X,Y )  = \Tr(XY)$.
Instead we write an element in the Cartan subalgebra $q = \sum_{i=1}^{r+1} q_i E_{ii}$ with the constraint $\sum q_i = 0$. 
In this description the simple roots are 
\begin{equation}
    \alpha_i(q) = q_i - q_{i+1}.
\end{equation}
The dynamical variables of the open Toda chain are the two elements $q = \sum_{i = 1}^{r+1} q_i E_{ii}$ and $p~=~ \sum_{i=1}^{r+1}p_iE_{ii}$ with the constraints $\sum_i q_i = 0$ and $\sum_i p_i = 0$. These constraints are not compatible with the canonical Poisson bracket, eq. (\ref{eq:Foundamental_Poisson}), so we use instead the Dirac bracket:
\begin{equation}
    \{ q_i, p_j \} = \frac{1}{2} \left( \delta_{ij} - \frac{1}{r+1} \right), \quad \{ p_i, p_j \} = \{ q_i, q_j \} = 0.
\end{equation}
The Hamiltonian 
\begin{equation}
    H = \sum_{i=1}^{r+1} p_i^2 + 2 \sum_{i=1}^r \mathrm{e}^{2(q_i - q_{i+1})}
\end{equation}
generates the equations of motion:
\begin{align}
    \dot{q}_i &= p_i \\
    \dot{p}_1 &= - 2 \mathrm{e}^{2(q_1 - q_2)} \\
    \dot{p}_i &= - 2 \mathrm{e}^{2(q_{i} - q_{i+1})} + 2 \mathrm{e}^{2(q_{i-1} - q_{i})} \\
    \dot{p}_{r+1} &= 2 \mathrm{e}^{2(q_r - q_{r+1})}.
\end{align}
The particular form of the equations for $\dot{p}_1$ and $\dot{p}_{r+1}$ explains the terminology \textit{open} Toda chain.

The equations of motion of the open Toda chain can be written in Lax form and the system is integrable.
Note that the exponential dependence in the Lax matrix is chosen so that the commutator $[M,L]$ produces terms proportional to $e^{2\alpha(q)}$, matching the potential term in the Hamiltonian.

\begin{proposition}
    The equations of motion of the open Toda chain admit a Lax pair representation $\dot{L} = {[M,L]}$ with 
    \begin{equation} \label{eq:lax_matrix_toda_chain}
        L = p + \sum_{\alpha \text{ simple}} n_\alpha \mathrm{e}^{\alpha(q)} (E_\alpha + E_{-\alpha}), \quad M = - \sum_{\alpha \text{ simple}} n_\alpha \mathrm{e}^{\alpha(q)} (E_\alpha - E_{-\alpha}).
    \end{equation}
\end{proposition}
\begin{proof}
   We compute the time derivative of $L$:
   \begin{equation}
       \dv{L}{t} = \dv{p}{t} + \sum_\alpha n_\alpha \alpha \left( \dv{q}{t} \right) \mathrm{e}^{\alpha(q)} (E_\alpha + E_{-\alpha}).
   \end{equation}
   Then we compute the commutator
   \begin{equation}
       [M,L] = \sum_\alpha n_\alpha \alpha(p) \mathrm{e}^{\alpha(q)} (E_\alpha + E_{- \alpha}) - \sum_\alpha \sum_\beta n_\alpha n_\beta \mathrm{e}^{(\alpha + \beta)(q)} [E_\alpha - E_{-\alpha}, E_\beta + E_{-\beta}].
   \end{equation}
   Since $[E_{\pm \alpha}, E_{\pm \beta}]$ is antisymmetric in the exchange of $\alpha$ and $\beta$, the second term reduces to 
   \begin{equation}
    - \sum_\alpha \sum_\beta n_\alpha n_\beta \mathrm{e}^{(\alpha + \beta)(q)} ([E_\alpha, E_{-\beta}] - [E_{-\alpha},E_{\beta}]).
   \end{equation}
   Since $\alpha$ and $\beta$ are simple roots, we have $[E_\alpha, E_{-\beta}] = \delta_{\alpha\beta}( E_{\alpha}, E_{-\alpha} ) H_\alpha$
   and therefore
   \begin{equation}
       [M,L] = -2\sum_\alpha \mathrm{e}^{2\alpha(q)} H_\alpha + \sum_{\alpha} n_\alpha \alpha(p) \mathrm{e}^{\alpha(q)} (E_\alpha + E_{-\alpha})
   \end{equation}
   which reproduces exactly the equations of motion, hence proving the Lax representation.
\end{proof}


\subsection[The \textit{r}-matrix of the Toda models]{The $\boldsymbol{r}$-matrix of the Toda models}\label{sec:toda_rmatrix}
The $r$-matrix of the Toda chain provides a canonical example of a constant $r$-matrix associated with a simple Lie algebra, and is known to satisfy the modified classical Yang–Baxter equation. In this section, we derive its explicit form and verify the associated Poisson structure.

\begin{theorem}
    Let $L$ be the Lax matrix $(\ref{eq:lax_matrix_toda_chain})$, with $p$ and $q$ satisfying the Poisson bracket relation $(\ref{eq:poison_bis})$.
    Then there exists an antisymmetric $r$-matrix 
    $r_{12} = - r_{21} \in \mathcal{G}^{\otimes 2}$,
    independent of the phase space variables $p$ and $q$, such that
    \begin{equation}
        \{ L_1 , L_2 \} = [ r_{12},\, L_1 + L_2 ].
    \end{equation}
    It is given by
    \begin{equation} \label{eq:r_12_definition}
        r_{12} = \frac12 \sum_{\alpha\ \text{positive}}
        \frac{E_{\alpha} \otimes E_{-\alpha} \;-\;
              E_{-\alpha} \otimes E_{\alpha}}
             {( E_{\alpha}, E_{-\alpha} )}.
    \end{equation}
    This formula implies that the conserved Hamiltonians are in involution.
\end{theorem}


The proof relies on the following auxiliary results.

\begin{lemma} \label{lemma:r_H}
    Let $r_{12}$ be defined by \eqref{eq:r_12_definition}.  
    Then
    \begin{equation}
        \left[r_{12},\,H\otimes1+1 \otimes H\right]=0
        \qquad \forall\, H \in \mathcal H.
    \end{equation}
\end{lemma}
\begin{proof}
    We directly compute the commutator:
    \begin{equation}
        [E_\alpha \otimes E_{-\alpha}, H_i \otimes 1 + 1 \otimes H_i] = [E_\alpha, H_i] \otimes E_{-\alpha} + E_{\alpha} \otimes [E_{-\alpha}, H_i].
    \end{equation}
    Using 
    \begin{equation}
        [H_i,E_\alpha] = \alpha(H_i) E_\alpha
        \qquad
        [H_i,E_{-\alpha}] = - \alpha(H_i) E_{-\alpha}
    \end{equation}
    we obtain
    \begin{equation}
        -\alpha(H_i) \, E_{\alpha} \otimes E_{-\alpha} + \alpha(H_i) \, E_{\alpha} \otimes E_{-\alpha} = 0.
    \end{equation}
    Similarly $ [E_{-\alpha}\otimes E_\alpha,\,H_i\otimes1+1 \otimes H_i] = 0 $. \\ 
    Thus every term in the sum commutes with $H_i\otimes1+1 \otimes H_i$.
\end{proof}

\begin{lemma} \label{lemma:r_E}
    Let $r_{12} \in \mathcal{G} \otimes \mathcal{G}$ be defined by equation~\eqref{eq:r_12_definition}. 
    Then, for every simple root $\alpha$, one has
    \begin{equation}
        [r_{12},\, E_{\pm \alpha} \otimes 1 + 1 \otimes E_{\pm \alpha}]
        =
        \frac{1}{2}
        \bigl(
        H_{\alpha} \otimes E_{\pm \alpha}
        -
        E_{\pm \alpha} \otimes H_{\alpha}
        \bigr).
    \end{equation}
\end{lemma}
\begin{proof}
    We compute $ [r_{12}, E_\alpha\otimes1 + 1 \otimes E_{\alpha}]$.
    For simplicity, we suppress normalization factors depending on the roots. A fully explicit computation requires keeping track of the coefficients $(E_\beta, E_{-\beta})$, but does not affect the structure of the result.
    \begin{align}
        &= \frac12 \sum_{\beta \text{ positive}}
        [E_\beta\otimes E_{-\beta}-E_{-\beta}\otimes E_\beta,\,
        E_\alpha \otimes 1 + 1 \otimes E_\alpha ] \\
        &= \frac{1}{2} \sum_{\beta \text{ positive}}
        [E_\beta,E_\alpha]\otimes E_{-\beta}
        + E_\beta \otimes [E_{-\beta},E_\alpha]
        - [E_{-\beta},E_\alpha]\otimes E_\beta
        - E_{-\beta}\otimes [E_\beta,E_\alpha].
    \end{align}
    We divide the sum into two contributions: $\alpha = \beta$ and $\alpha \neq \beta$.
    Only the contribution with $\alpha=\beta$ produces Cartan elements, since $[E_\alpha,E_{-\alpha}] = H_\alpha$. We obtain:
    \begin{equation}
        \frac12
        \left(
        H_\alpha \otimes E_\alpha - E_\alpha \otimes H_\alpha
        \right).
    \end{equation}
    For every $\beta\neq\alpha$, the commutator $[r_{12}, E_{\pm\alpha} \otimes 1 + 1 \otimes E_{\pm\alpha}]$ is a linear combination of tensors of the form
    \begin{equation}
        E_{\beta\pm\alpha}\otimes E_{-\beta}, \qquad
        E_{\pm\alpha-\beta}\otimes E_{\beta}
    \end{equation}
    whenever these root spaces are nonzero. \\
    Consider the first type of terms: $E_{\beta + \alpha} \otimes E_{-\beta}$. We have
    \begin{align}
        [E_\beta \otimes E_{-\beta},\, E_\alpha \otimes 1] &= c_1 \, E_{\beta + \alpha} \otimes E_{-\beta} \\
        [E_{\alpha + \beta} \otimes E_{-\beta - \alpha},\, 1 \otimes E_{\alpha}] &= c_2 \, E_{\beta + \alpha} \otimes E_{-\beta}.
    \end{align}
    Hence,
    \begin{equation}
        [E_\beta \otimes E_{-\beta},\, E_\alpha \otimes 1] 
        + [E_{\alpha + \beta} \otimes E_{-\beta - \alpha},\, 1 \otimes E_{\alpha}]
        = (c_1 + c_2)\, E_{\beta + \alpha} \otimes E_{-\beta}.
    \end{equation}
   Using the invariance of the scalar product, i.e.
    \begin{equation}
        (X,[Y,Z]) = ([X,Y],Z),
    \end{equation}
    we compute
    \begin{align}
        c_1 &= (E_{-\beta-\alpha},[E_\beta,E_\alpha]), \\
        c_2 &= (E_\beta,[E_{-\beta-\alpha},E_\alpha])
            = -([E_\beta,E_\alpha],E_{-\beta-\alpha})
            = -c_1.
    \end{align}
    Therefore, the contributions of this type cancel, and their sum vanishes.
    Consider now the second type: $ E_{\alpha - \beta} \otimes E_\beta $.\\
    Since $\alpha$ is a simple root, whenever $\beta - \alpha$ is a root we may set $\gamma = \beta - \alpha$, so that $\beta = \alpha + \gamma$.
    Then
    \begin{equation}
        E_{\alpha - \beta} \otimes E_\beta = E_{-\gamma} \otimes E_{\alpha + \gamma}.
    \end{equation}
    This reduces the expression to the same form treated in the previous case, so the argument proceeds identically.

    Therefore, the only remaining terms are those with $\alpha = \beta$, which proves the claim:
    \begin{equation}
        [r_{12},\, E_{\pm \alpha} \otimes 1 + 1 \otimes E_{\pm \alpha}]
        = \frac{1}{2}\bigl(H_{\alpha} \otimes E_{\pm \alpha} - E_{\pm \alpha} \otimes H_{\alpha}\bigr).
    \end{equation}
\end{proof}
We now return to the proof of the theorem.
\begin{proof}[Proof of the theorem]
    We first compute the Poisson bracket of the Lax matrix:
    \begin{equation}
        \{L_1,L_2\}
        =
        \frac12 \sum_{\alpha\ \text{simple}}
        n_\alpha \mathrm{e}^{\alpha(q)}
        \bigl[
        H_\alpha \otimes (E_\alpha + E_{-\alpha})
        -
        (E_\alpha + E_{-\alpha}) \otimes H_\alpha
        \bigr].
    \end{equation}

    We now show that the tensor $r_{12}$ defined in \eqref{eq:r_12_definition} reproduces the same expression through the commutator $[r_{12},L_1+L_2]$.

    Since
    \begin{equation}
        L
        =
        p
        +
        \sum_{\alpha\ \text{simple}}
        n_\alpha \mathrm{e}^{\alpha(q)} (E_\alpha+E_{-\alpha}),
    \end{equation}
    we have
    \begin{equation}
        L_1+L_2
        =
        p \otimes 1 + 1 \otimes p
        +
        \sum_{\alpha\ \text{simple}}
        n_\alpha \mathrm{e}^{\alpha(q)}
        \bigl(
        (E_\alpha+E_{-\alpha}) \otimes 1
        +
        1 \otimes (E_\alpha+E_{-\alpha})
        \bigr).
    \end{equation}
    Therefore
    \begin{align}
        [r_{12},L_1+L_2]
        &=
        [r_{12},p \otimes 1 + 1 \otimes p] \notag \\
        &\quad +
        \sum_{\alpha\ \text{simple}}
        n_\alpha \mathrm{e}^{\alpha(q)}
        [r_{12},(E_\alpha+E_{-\alpha}) \otimes 1 + 1 \otimes (E_\alpha+E_{-\alpha})].
    \end{align}
    Since $p=\sum_i p_i H_i$, Lemma~\ref{lemma:r_H} implies
    \begin{equation}
        [r_{12},p \otimes 1 + 1 \otimes p]
        =
        \sum_i p_i [r_{12},H_i \otimes 1 + 1 \otimes H_i]
        =
        0.
    \end{equation}
    For the second term, by linearity of the commutator and by Lemma~\ref{lemma:r_E}, valid for every simple root $\alpha$, we obtain
    \begin{align}
        & [r_{12},(E_\alpha+E_{-\alpha}) \otimes 1 + 1 \otimes (E_\alpha+E_{-\alpha})] \notag \\
        &= [r_{12},E_\alpha \otimes 1 + 1 \otimes E_\alpha]
        +
        [r_{12},E_{-\alpha} \otimes 1 + 1 \otimes E_{-\alpha}] \notag \\
        &= \frac12
        \bigl(
        H_\alpha \otimes E_\alpha - E_\alpha \otimes H_\alpha
        \bigr)
        +
        \frac12
        \bigl(
        H_\alpha \otimes E_{-\alpha} - E_{-\alpha} \otimes H_\alpha
        \bigr) \notag \\
        &= \frac12
        \bigl[
        H_\alpha \otimes (E_\alpha+E_{-\alpha})
        -
        (E_\alpha+E_{-\alpha}) \otimes H_\alpha
        \bigr].
    \end{align}

    Substituting these identities into the expression of $[r_{12},L_1+L_2]$, we find
    \begin{equation}
        [r_{12},L_1+L_2]
        =
        \frac12 \sum_{\alpha\ \text{simple}}
        n_\alpha \mathrm{e}^{\alpha(q)}
        \bigl[
        H_\alpha \otimes (E_\alpha + E_{-\alpha})
        -
        (E_\alpha + E_{-\alpha}) \otimes H_\alpha
        \bigr].
    \end{equation}

    This coincides exactly with the expression of $\{L_1,L_2\}$. Hence
    \begin{equation}
        \{L_1,L_2\} = [r_{12},L_1+L_2].
    \end{equation}
    This proves the theorem.
\end{proof}


An alternative and more structural formulation of the $r$-matrix formalism can be obtained by introducing a linear operator $R:\mathcal G \to \mathcal G$ associated with $r_{12}$.
This point of view, which is closely related to the classical Yang--Baxter equation, is presented in Appendix~\ref{app:Rmatrix}.

\section{Independence of \texorpdfstring{${H_k}$}{Hk}}\label{sec:Independence of Hk}
In this section we prove the independence of the Hamiltonians associated with the Lax matrix.

\begin{proposition}
    In the case $\mathfrak{sl}(r+1)$, the functions $H_2,\dots,H_{r+1}$, with
    \begin{equation}
        H_k=\Tr\left(L^k\right),
    \end{equation}
    are functionally independent on a nonempty open subset of phase space.
\end{proposition}

\begin{proof}
    In the case $\mathfrak{sl}(r+1)$ one has $\Tr L = 0$, hence $H_1=0$ identically.
    Therefore we restrict to the functions $H_2,\dots,H_{r+1}$, which provide $r$ functions.

    The variables $p_i$ satisfy the constraint $\sum_{i=1}^{r+1} p_i = 0$, so we may regard $p_1,\dots,p_r$ as independent variables.

    For the Toda Lax matrix, one has
    \begin{equation}
        H_k = \sum_{i=1}^{r+1} p_i^k + R_k(p,q),
    \end{equation}
    where $R_k(p,q)$ has degree in $p$ strictly smaller than $k$.

    We first prove that the functions
    \begin{equation}
        F_k(p_1,\dots,p_r) = \sum_{i=1}^{r+1} p_i^k,
        \qquad k=2,\dots,r+1,
    \end{equation}
    are functionally independent on the hyperplane $\sum_i p_i = 0$.

    Writing
    \begin{equation}
        p_{r+1} = -\sum_{i=1}^r p_i,
    \end{equation}
    we compute the Jacobian of the functions $F_k$ with respect to the independent variables $p_1,\dots,p_r$:
    \begin{equation}
        \pdv{F_k}{p_j}
        = k p_j^{k-1} + k p_{r+1}^{\,k-1}\pdv{p_{r+1}}{p_j}
        = k\left(p_j^{k-1} - p_{r+1}^{\,k-1}\right),
        \qquad j=1,\dots,r.
    \end{equation}

    Thus the Jacobian matrix $J_F$ is
    \begin{equation}
        (J_F)_{kj} = k\left(p_j^{k-1} - p_{r+1}^{\,k-1}\right),
        \qquad k=2,\dots,r+1,\quad j=1,\dots,r.
    \end{equation}
    
    Up to the nonzero factors $k$, this matrix has the structure of a Vandermonde-type matrix built from the variables $p_1,\dots,p_{r+1}$.
    A direct computation shows that its determinant is proportional to
    \begin{equation}
        \prod_{1 \le i < j \le r+1} (p_i - p_j),
    \end{equation}
    and is therefore nonzero whenever the $p_i$ are pairwise distinct.
    Hence the Jacobian $J_F$ is non-degenerate on a nonempty open subset of the hyperplane $\sum_i p_i=0$, and the functions $F_2,\dots,F_{r+1}$ are functionally independent there.

    We now return to the functions $H_k$. Their Jacobian satisfies
    \begin{equation}
        J_H = J_F + J_R,
        \qquad
        (J_R)_{kj} = \pdv{R_k(p,q)}{p_j}.
    \end{equation}
    Since $\deg_p R_k < k$, we have
    \begin{equation}
        \deg_p \pdv{R_k}{p_j} < k-1,
    \end{equation}
    so every entry of $J_R$ has degree in $p$ strictly smaller than the corresponding entry of $J_F$.

    It follows that $\det J_H$ has the same leading homogeneous term, as a polynomial in $p$, as $\det J_F$.
    In particular, since $\det J_F$ is not identically zero, neither is $\det J_H$.

    Hence $\det J_H \neq 0$ on a nonempty open subset of phase space.
    Therefore the functions $H_2,\dots,H_{r+1}$ are functionally independent on a nonempty open subset of phase space.
\end{proof}


\section{Conclusions}
In this work we have developed the $r$-matrix formulation of integrable systems starting from the Lax pair description and progressively uncovering the underlying algebraic structure.

We first recalled that the existence of a Lax representation allows one to construct conserved quantities as spectral invariants of the Lax matrix. We then introduced the Poisson algebra of Lax matrices and showed that, when the Poisson brackets take the $r$-matrix form
\begin{equation}
    \{L_1,L_2\} = [r_{12},L_1] - [r_{21},L_2],
\end{equation}
the spectral invariants $\Tr(L^n)$ are in involution. We also discussed the converse statement, namely the local reconstruction of an $r$-matrix structure from the involution of the eigenvalues, under suitable regularity assumptions.

The algebraic content of this formalism was then developed through the tensor Casimir element and the classical and modified Yang--Baxter equations. In particular, we showed how the Yang--Baxter equation appears as the compatibility condition ensuring the Jacobi identity for the Poisson bracket, and how the Casimir element naturally enters the structure of the theory.

These general constructions were finally illustrated on the example of the open Toda chain associated with a simple Lie algebra. We reviewed its Lax representation and derived explicitly the corresponding $r$-matrix, showing that it reproduces the Poisson structure and yields commuting Hamiltonians. This provides a concrete realization of the abstract framework and clarifies how the general algebraic formalism is implemented in an integrable model.

For completeness, we also briefly reformulated the tensorial picture in terms of the corresponding $R$-matrix operator in the appendix.

Overall, the $r$-matrix formalism provides a unified perspective on integrability, linking Hamiltonian dynamics, spectral invariants, and Lie algebraic structures in a natural and explicit way.

\newpage
\appendix
\section{\texorpdfstring{$R$}{R}-matrix formulation}\label{app:Rmatrix}
In this appendix we reformulate the $r$-matrix structure in operator form \cite{SemenovTyanShanskii1983,babelon2003introduction}.

We fix a matrix realization $\mathcal G \subset \mathfrak{gl}(r+1)$ and use the invariant bilinear form
\begin{equation}
    (X,Y)=\Tr(XY),
\end{equation}
which identifies $\mathcal{G} \simeq \mathcal{G}^*$.

\begin{definition}
    For $L \in \mathcal{G}$ and $X \in \mathcal{G}$, we define the scalar function
    \begin{equation}
        L(X) \equiv (L,X).
    \end{equation}
\end{definition}

Using this identification, any element $r_{12} \in \mathcal{G}^{\otimes 2}$ defines a linear operator
$R:\mathcal{G} \to \mathcal{G}$ by
\begin{equation}
    R(X)=\Tr_2\left(r_{12}(1 \otimes X)\right).
\end{equation}
This provides an equivalent operator formulation of the $r$-matrix formalism.

\begin{proposition}
    Let $r_{12} \in \mathcal{G}^{\otimes 2}$ and let $R$ be the associated operator.  
    If the Poisson brackets of $L \in \mathcal{G}$ satisfy
    \begin{equation}
        \{L_1,L_2\} = [r_{12},L_1] - [r_{21},L_2],
    \end{equation}
    then, for all $X,Y \in \mathcal{G}$,
    \begin{equation}
        \{L(X),L(Y)\} = L([X,Y]_R),
    \end{equation}
    where
    \begin{equation}
        [X,Y]_R = [R(X),Y] + [X,R(Y)].
    \end{equation}
\end{proposition}

\begin{proof}[Sketch of proof]
    Pairing the tensorial identity with $X \otimes Y$ and using the invariance of the scalar product,
    one rewrites both sides in terms of the pairing $(L,X)$. The commutator terms reduce to expressions involving $R$, yielding the stated formula.
\end{proof}

The Yang--Baxter equation admits a natural formulation in terms of the operator $R$.

\begin{proposition}
    Let $r_{12} \in \mathcal{G}^{\otimes 2}$ be antisymmetric and let $R$ be defined as above.  
    If $r_{12}$ satisfies the modified classical Yang--Baxter equation
    \begin{equation}
        [[r,r]] = \mu \Omega,
    \end{equation}
    then $R$ satisfies
    \begin{equation}
        [R(X),R(Y)] - R\left([R(X),Y] + [X,R(Y)]\right) = -\mu [X,Y]
        \qquad \forall X,Y \in \mathcal{G}.
    \end{equation}
\end{proposition}

\begin{proof}[Sketch of proof]
    Pairing the tensorial identity $[[r,r]]=\mu\Omega$ with $X \otimes Y \otimes Z$
    and using invariance, one translates the tensorial commutators into operator expressions involving $R$.
\end{proof}

In the case of the Toda chain, the $r$-matrix constructed in Section~\ref{sec:toda_rmatrix}
defines an operator $R$ diagonal in the root space decomposition:
\begin{equation}
    R(H)=0, \qquad
    R(E_{\pm\beta}) = \pm \frac12 E_{\pm\beta}
    \qquad (\beta \text{ positive}).
\end{equation}
This provides an explicit solution of the modified classical Yang--Baxter equation. 

\newpage
\nocite{*}
\printbibliography

\end{document}
